% ============================================================
% CHAPTER 4 — System Design and Methodology (REVISED v2)
% ============================================================

\chapter{System Design and Methodology}
\label{chap:method}


As discussed in Chapter~\ref{chap:literarture}, existing \ac{RL}
approaches to process optimisation typically operate over imperative
models or unconstrained action spaces, so invalid traces can be
generated during training~\cite{Huang2011ReinforcementManagement,
Kubrak2021PrescriptiveVadis}, and few works combine declarative
semantics, multi-objective optimisation, and resource-aware decisions
in one framework~\cite{Diaz2024Pareto-OptimalModels,
KirchdorferLukasFromPolicies, Odriozola-Olalde2025TowardsFramework}.


This chapter presents the proposed framework in four parts.
Section~\ref{sec:compliance} integrates \ac{DCR} semantics into a
\ac{RL} environment with a shield, intrinsic to the execution engine,
that guarantees compliance-by-construction during
training~\cite{Alshiekh2017SafeShielding} (contribution~C1).
Section~\ref{sec:reward} introduces a normalised
multi-objective reward that scalarises cost and duration under
user-defined weights~\cite{Roijers2013ADecision-Making}, and
Section~\ref{sec:mooptimisation} explains how this approximates the
Pareto front (contribution~C2). Section~\ref{sec:roles} extends the
environment to joint \textit{(event, role)} decisions
(contribution~C3). Section~\ref{sec:ppo} describes the \ac{PPO}
training setup.

% ============================================================
\section{System Architecture}
\label{sec:architecture}

The framework connects a declarative \ac{DCR} model with a \ac{PPO}
agent in a closed training loop (Figure~\ref{fig:architecture}). The
design separates \emph{process logic}, which actions are valid and
when, from \emph{learning logic}, which actions are preferable and
why. In Figure~\ref{fig:architecture}, process logic is the
responsibility of the \ac{DCR} execution engine; learning logic is the
responsibility of the \ac{PPO} agent and the reward function; the
environment in between only translates one representation into the
other and makes neither kind of decision itself. Only the \ac{DCR}
execution engine can modify the process state, and the shield that
protects that state is part of the engine itself.


\begin{figure}[H]
\centering
\includegraphics[width=\textwidth]{Pictures/thesis_architecture.png}
\caption{Overall system architecture. The \ac{PPO} agent (left)
communicates with the \ac{DCR} execution engine (right) through the
Node Adapter \ac{RL} environment (centre). The shield is part of the
engine: any proposed action whose event is not enabled is intercepted
before execution, so no invalid transition can modify the process
state. Dashed arrows indicate logging and outputs; solid arrows
indicate data flow within one environment step.}
\label{fig:architecture}
\end{figure}

\textbf{Input.} A \ac{DCR} graph in XML, with activities, relations,
cost and duration annotations, and optional role profiles, plus an
experiment configuration with the weights $(\alpha, \beta)$, \ac{PPO}
hyperparameters, and role multipliers.

\textbf{\ac{PPO} agent.}
At each step the agent observes the encoded \ac{DCR} marking and
selects an action from the full discrete action space; no mask
restricts this choice, so the agent may propose actions that are not
currently enabled. The algorithm, network architecture, and
hyperparameters are detailed in Section~\ref{sec:ppo}.


\textbf{\ac{DCR} \ac{RL} environment.}
The translation layer between the Python agent and the TypeScript
engine. Its \emph{Action Decoder} maps the selected index to an
\texttt{(event, role)} pair; its \emph{State Encoder} converts the
updated marking into the observation vector. The environment also
logs the engine's enabled-set as diagnostic information, but never
uses it to restrict the agent (a deliberate decision, see
Section~\ref{sec:shielding}).

\textbf{\ac{DCR} execution engine.}
Implements the formal \ac{DCR}
semantics~\cite{Diaz2024Pareto-OptimalModels}: it tracks the marking,
computes the enabled events, updates the state after execution, and
detects acceptance. The shield is an inherent property of this
engine. If the proposed event is not enabled, the engine returns
$r = -10$ and leaves the marking unchanged
(Section~\ref{sec:shielding}).

\textbf{Output.} Each run produces the trained policy for its
$(\alpha, \beta)$ configuration, the execution traces generated under
that policy with their cost and duration, and training statistics
(reward, illegal-action ratio, acceptance rate, role selections) used
to build the Pareto front of Chapter~\ref{chap:results}.

% ============================================================
\section{Compliance-by-Construction via DCR Shielding}
\label{sec:compliance}

This section formalises contribution~C1. Standard \ac{RL} agents
explore freely and may find reward-maximising behaviour that violates
process constraints. Penalty-based and Lagrangian
approaches~\cite{AltmanCONSTRAINEDPROCESSES, Achiam2017ConstrainedOptimization,
Huang2011ReinforcementManagement} tolerate these violations during
training and only discourage them statistically. Here, instead,
violations are structurally impossible from the first
episode~\cite{Alshiekh2017SafeShielding}: the engine intercepts any non-enabled
action before the marking can change, so every trace generated during
training consists only of valid \ac{DCR} transitions, regardless of
reward design or training budget.

\subsection{State Representation}
\label{sec:state}

The agent observes the marking
$M = (\mathit{In}, \mathit{Ex}, \mathit{Pe})$ as a binary vector.
Each event $e \in \mathcal{E}$ contributes three
features~\cite{Diaz2024Pareto-OptimalModels, Trinh2024OnOfConstraints}:

\begin{equation}
\mathbf{o}_e =
\bigl[\mathbf{1}[e \in \mathit{In}],\;
      \mathbf{1}[e \in \mathit{Ex}],\;
      \mathbf{1}[e \in \mathit{Pe}]\bigr]
\in \{0,1\}^3
\end{equation}

and the observation is their concatenation,
$\mathbf{o} \in \{0,1\}^{3|\mathcal{E}|}$. The dimensions per
benchmark are given in Chapter~\ref{chap:experimental_setup}.

\subsection{Action Space}
\label{sec:action_space}

The base action space has one action per event. The role-conditioned
extension (Section~\ref{sec:roles}) expands it to one action per
\texttt{(event, role)} pair, giving
$|\mathcal{A}| = |\mathcal{E}| \times |\mathcal{R}|$. The agent
samples from this full space at every step, including non-enabled
events; whether the action is executed is decided by the shield.

\subsection{Enabled Events and Accepting States}
\label{sec:dcr_semantics}

Following Definition~\ref{def:dcr_transitions}, an event is
\emph{enabled} if it is included, all its conditions have been
executed, and no blocking milestones remain pending. A marking is
\emph{accepting} when no event is both included and pending:

\begin{equation}
\mathit{In} \cap \mathit{Pe} = \emptyset
\label{eq:accepting}
\end{equation}

Accepting states correspond to valid process completions.

\subsection{Shielding Mechanism}
\label{sec:shielding}

Figure~\ref{fig:shielding} shows the mechanism. At every step, the
engine checks the proposed action against the enabled set
\emph{before} executing it. If the action is enabled, it is executed
and the marking advances. If not, the action is intercepted, the
marking stays exactly as it was, and the agent receives $r = -10$
with the episode continuing.

In the placement taxonomy of Alshiekh et
al.~\cite{Alshiekh2017SafeShielding}, the proposed action is examined only
after the agent has selected it, which makes this a \emph{post-posed} shield
rather than a \emph{preemptive} one; a preemptive shield would instead restrict
the action set offered to the agent in advance. It departs from the standard
post-posed formulation in one respect: a non-enabled action is rejected
outright and the marking is left unchanged, rather than being overwritten by a
safe substitute.

Two designs could enforce this guarantee: masking the action space so
the agent can only choose enabled actions, or letting the agent
choose freely and intercepting invalid choices. The framework
deliberately uses interception. The guarantee is the same, no invalid
transition ever reaches the marking, but the agent's illegal attempts
remain observable: each blocked attempt gives a learning signal
through the penalty, and the illegal-attempt rate over training
measures how well the agent has learned the legal structure of the
process. Chapter~\ref{chap:internal_validation} uses exactly this
property.

Table~\ref{tab:shield-trace} shows the mechanism on real data: the
first training episode of the Loan Application experiment, when the
agent still acts randomly, a worst case for the shield. The 20
intercepted attempts never change $|\mathit{Pending}|$; the 9 legal
transitions advance the marking until acceptance at step 29.

\begin{table}[H]
\centering
\footnotesize
\caption{First training episode of the Loan Application roles experiment
($\alpha{=}0$, $\beta{=}0$, seed $s{=}1$): 9 legal transitions, 20
intercepted attempts, acceptance at step 29. Intercepted attempts leave
$|\mathit{Pending}|$ unchanged; the full 29-step trace is given in
Appendix~\ref{app:shield_trace}.}
\label{tab:shield-trace}
{\renewcommand{\arraystretch}{1.3}
\setlength{\tabcolsep}{4pt}
\begin{tabular}{r l p{4.3cm} c r}
\toprule
\textbf{Step} & \textbf{Type} & \textbf{Activity} &
\textbf{$|\mathit{Pending}|$} & \textbf{Reward} \\
\midrule
1  & \textcolor{BrickRed}{\textsf{illegal\;\ding{55}}} & Cancel application & $0 \!\to\! 0$ & $-10$  \\
2  & legal\;\ding{51} & Check application form completeness & $0 \!\to\! 3$ & $-0.5$ \\
3  & \textcolor{BrickRed}{\textsf{illegal\;\ding{55}}} & Cancel application & $3 \!\to\! 3$ & $-10$  \\
4  & legal\;\ding{51} & Appraise property & $3 \!\to\! 3$ & $-0.5$ \\
5  & \textcolor{BrickRed}{\textsf{illegal\;\ding{55}}} & Design loan offer & $3 \!\to\! 3$ & $-10$  \\
6  & legal\;\ding{51} & Check credit history & $3 \!\to\! 2$ & $+1.5$ \\
7  & \textcolor{BrickRed}{\textsf{illegal\;\ding{55}}} & Approve loan offer & $2 \!\to\! 2$ & $-10$  \\
8  & legal\;\ding{51} & Appraise property (repeated) & $2 \!\to\! 2$ & $-1.5$ \\
9  & \textcolor{BrickRed}{\textsf{illegal\;\ding{55}}} & Approve application & $2 \!\to\! 2$ & $-10$  \\
10 & legal\;\ding{51} & AML check & $2 \!\to\! 1$ & $+1.5$ \\
11 & legal\;\ding{51} & Appraise property (repeated) & $1 \!\to\! 1$ & $-1.5$ \\
\multicolumn{5}{c}{\footnotesize\textit{$\cdots$\;steps 12--28: 15 further intercepted attempts and 2 legal repetitions, $|\mathit{Pending}|=1$ throughout\;$\cdots$}} \\
29 & \textbf{accept\;\ding{51}} & Assess loan risk & $1 \!\to\! 0$ & $+100$ \\
\bottomrule
\end{tabular}}
\end{table}

The $-10$ penalty plays no role in the compliance guarantee, which
comes from the environment design, not from the penalty
size~\cite{Odriozola-Olalde2025TowardsFramework}; it only speeds up
learning. The guarantee should also be stated precisely. What the
shield makes logically necessary is the \emph{validity} of every
trace: no trace can contain an invalid transition. Reaching an
accepting state is not guaranteed; an episode may run out of steps
and end as non-accepting (Section~\ref{sec:training_pipeline}). The
acceptance rates in Chapter~\ref{chap:results} are therefore learned,
empirical results, while the validity of every trace is structural.

% ============================================================
\section{Reward Design}
\label{sec:reward}

This section formalises the first part of contribution~C2.
Compliance is already guaranteed by construction, so the reward's job
is to distinguish between valid traces: a purely structural reward
would leave the agent indifferent to cost and duration, converging to
compliant but suboptimal executions.
Table~\ref{tab:reward_components} lists every component of the
per-step signal.

\begin{table}[H]
\centering
\caption{Reward components. Structural components and the
cost--duration penalties apply to legal, non-terminal actions; the
illegal penalty and the terminal bonus each replace the whole step
reward when triggered.}
\label{tab:reward_components}
{\renewcommand{\arraystretch}{1.5}
\setlength{\tabcolsep}{14pt}
\begin{tabular}{l l p{4.6cm}}
\toprule
\textbf{Component} & \textbf{Value} & \textbf{Trigger} \\
\midrule
Illegal action penalty & $-10$ &
  Proposed event not enabled; shield intercepts \\
Terminal bonus & $+100$ &
  Accepting state reached \\
Base step reward & $+1$ &
  Every legal, non-terminal action \\
Step penalty & $-1.5$ &
  Every legal, non-terminal action; discourages long episodes \\
Progress bonus & $+2$ per event &
  Each pending event resolved for the first time in the episode \\
Repetition penalty & $-1$ &
  Event already executed earlier in the episode (any role) \\
Cost penalty & $-\alpha \cdot c_{e,\rho} / c_{\max}$ &
  Each legal action; role-adjusted cost \\
Duration penalty & $-\beta \cdot d_{e,\rho} / d_{\max}$ &
  Each legal action; role-adjusted duration \\
\bottomrule
\end{tabular}}
\end{table}


The complete per-step reward is:

\begin{equation}
r_t =
\begin{cases}
-10 & \text{if } a \notin \mathit{enabled}(M) \\[4pt]
+100 & \text{if the resulting marking is accepting} \\[4pt]
1 - 1.5
+ 2\,n_{\text{res}}
- \mathbf{1}[\text{repeat}]
- \alpha\dfrac{c_{e,\rho}}{c_{\max}}
- \beta\dfrac{d_{e,\rho}}{d_{\max}}
& \text{otherwise}
\end{cases}
\label{eq:reward_total}
\end{equation}

where $n_{\text{res}}$ is the number of pending events resolved for
the first time in this step and $\mathbf{1}[\text{repeat}]$ is 1 if
the event was already executed in the episode. The step penalty is
configurable and was $-1.5$ in all experiments.

The auxiliary shaping terms, a per-step penalty discouraging long
episodes, a progress bonus, and a repetition penalty---follow standard
reward-design practice for episodic \ac{RL}~\cite{Sutton2018ReinforcementIntroduction,
Mataric1994RewardLearning}, where Mataric's progress estimators are a
direct precedent for rewarding incremental advancement. The
potential-based analysis of Ng et al.~\cite{NgPolicyShaping}
shows when such terms preserve the optimal policy; our progress and
repetition terms are history-dependent and therefore not
potential-based, of the form $\gamma\Phi(s')-\Phi(s)$, so they carry
no such guarantee, they are deliberate biases, most visibly the
brevity bias of Section~\ref{sec:discussion_reward}. The magnitudes
were fixed empirically during development, not derived, and not
tuned by systematic search.


\textbf{Structural components.}
The terminal bonus of $+100$ is the main objective signal and is
large enough to dominate all other components, so completing the
process is always the preferred outcome. The base reward $+1$ keeps
the signal dense from the start of training, while the step penalty
$-1.5$ makes each extra step costly. Two small shaping terms refine
this: the progress bonus rewards resolving pending events, but only
the first time each event is resolved within an episode, which
prevents the agent from farming the bonus through cycles that resolve
and re-trigger the same obligation; the repetition penalty
discourages re-executing events, counted per event regardless of
role. Both are deliberate, lightweight reward shaping, small relative
to the terminal bonus, so they speed up convergence without changing
which traces are optimal. The $-10$ penalty, as noted above, is only
a learning signal.

\textbf{Normalised cost--duration penalty.}
\label{sec:reward_normalised}
Subtracting raw cost and duration values would create scale problems:
large durations could dominate the signal. Both terms are therefore
divided by the maximum base values in the graph, $c_{\max}$ and
$d_{\max}$, computed at load time, so all components stay on a
comparable scale. When roles are active, the multipliers of
Section~\ref{sec:roles} are applied to the base values \emph{before}
normalisation, so the normalised terms range over
$[0, \mu^\rho_{\max}]$ rather than $[0,1]$.

The three reward levels visible among the legal steps of
Table~\ref{tab:shield-trace} (where $\alpha = \beta = 0$) follow
directly from Equation~\ref{eq:reward_total}: no progress gives
$1 - 1.5 = -0.5$ (steps 2, 4), resolving one pending event gives
$1 - 1.5 + 2 = +1.5$ (steps 6, 10), and a repetition gives
$1 - 1.5 - 1 = -1.5$ (steps 8, 11).

% ============================================================
\section{Multi-objective Optimisation Strategy}
\label{sec:mooptimisation}

This section formalises the second part of contribution~C2: how the
reward of Section~\ref{sec:reward} is used to approximate the Pareto
front of compliant executions via linear scalarisation.

\paragraph{Scalarisation and the role of $(\alpha, \beta)$.}
\label{sec:scalarisation}
Although DCR execution is not a general stochastic MDP, each episode is a finite deterministic decision problem in which the next marking depends only on the current marking and the executed action, so the Markov property holds. Concretely, the MDP of Definition~\ref{def:mdp} is instantiated directly on top of the DCR labelled transition system of Definition~\ref{def:dcr_transitions}: states are markings, actions are events, and the transition function $P$ is the (deterministic) transition relation of the LTS itself, with $P(s'\mid s,a)=1$ for the unique $s'$ it produces. The LTS therefore supplies the dynamics; the reward function of Section~\ref{sec:reward} is the only addition needed to
frame the problem as an MDP.

The framework therefore falls within the multi-objective MDPs of Roijers et
al.~\cite{Roijers2013ADecision-Making}, and the scalarisation theory of Section~\ref{sec:lit_scalarisation} applies directly: training independent agents under different $(\alpha, \beta)$ weights and pooling their accepting traces approximates the Pareto front, with the known limitation that only solutions on the convex hull are
recoverable~\cite{Vamplew2008OnFronts}. For the benchmarks evaluated
here this does not affect the results: as Chapter~\ref{chap:results}
shows, the scalarisation recovers all known Pareto-optimal solutions,
consistent with a locally convex front.

The weight configurations follow the principle of axis-aligned
extremes, a symmetric centre, and asymmetric
probes~\cite{DasDr.Das-NBIPaper-1998}:

\begin{itemize}
  \item $(1,0)$ and $(0,1)$: extremes that isolate each objective and
        anchor the endpoints of the front;
  \item $(0.5,\ 0.5)$: the balanced configuration;
  \item $(2,\ 0.5)$ and $(0.5,\ 2)$: asymmetric probes of
        intermediate regions;
  \item $(0,\ 0)$: a baseline with no cost or duration incentive,
        used for convergence analysis and excluded from the Pareto
        front.
\end{itemize}

The values $(2, 0.5)$ are used instead of normalised equivalents like
$(0.8, 0.2)$ because, with both objectives already normalised, only
the ratio of $\alpha$ to $\beta$ matters; integer-like values are
simply easier to read. The mapping from weights to front points is
many-to-one over a discrete trace space, examined empirically in
Chapter~\ref{chap:internal_validation}.


% ============================================================
\section{PPO Training Setup}
\label{sec:ppo}

This section describes the concrete training setup that instantiates
the framework of Sections~\ref{sec:compliance}--\ref{sec:roles}: the
learning algorithm, its hyperparameters, and the training pipeline
shared by every experiment in this thesis.

\subsection{Algorithm}
\label{sec:ppo_algorithm}


The agent is trained with the standard \texttt{PPO} class of Stable
Baselines~3 (multilayer perceptron
policy)~\cite{Raffin2021Stable-Baselines3:Implementations}. \ac{PPO} fits this framework for
three reasons. First, the action space is discrete and small, between
5 and 24 actions depending on the benchmark, a regime where on-policy
policy-gradient methods work reliably. Second, on-policy learning
pairs naturally with the shield: every state visited in a rollout is
a valid \ac{DCR} marking, so the agent learns only over the legal
state space of the process. Third, Stable Baselines~3 integrates
cleanly with the Gymnasium environment~\cite{Towers2025Gymnasium:Environments}, avoiding a
custom training loop. The dense structural reward of
Section~\ref{sec:reward} reduces the credit assignment problem of a
purely terminal formulation, making training feasible within a moderate budget and without curriculum learning.
The policy network follows the Stable Baselines~3 \texttt{MlpPolicy}
default: two separate networks for the actor and the critic, each
with two hidden layers of 64 units and $\tanh$ activations, of
identical shape but with independent weights ($\mathtt{net\_arch} = \{\texttt{pi}: [64,64], \texttt{vf}: [64,64]\}$).



\subsection{Hyperparameters}
\label{sec:training_config}

Table~\ref{tab:hyperparams} lists the hyperparameters, shared by all
experiments. They follow the Stable Baselines~3 defaults, which
mirror Schulman et al.~\cite{Schulman2017ProximalAlgorithms}, except
the entropy coefficient, raised to $0.1$ to sustain exploration. No
hyperparameter search was performed; this is a threat to internal
validity discussed in Section~\ref{sec:limitations}.

\begin{table}[H]
\centering
\caption{\ac{PPO} hyperparameters, shared by all experiments.}
\label{tab:hyperparams}
{\renewcommand{\arraystretch}{1.5}
\setlength{\tabcolsep}{14pt}
\begin{tabular}{ll}
\toprule
\textbf{Hyperparameter} & \textbf{Value} \\
\midrule
Algorithm               & PPO (Stable Baselines 3, MlpPolicy) \\
Learning rate           & $3 \times 10^{-4}$ \\
Discount factor $\gamma$ & $0.99$ \\
Steps per update        & $2{,}048$ \\
Batch size              & $64$ \\
Epochs per update       & $10$ \\
Clip range              & $0.2$ \\
Entropy coefficient     & $0.1$ \\
Policy network          & MLP, 2 $\times$ 64 units, $\tanh$ \\
\bottomrule
\end{tabular}}
\end{table}


\subsection{Training Pipeline}
\label{sec:training_pipeline}

One independent agent is trained per scalarisation configuration.
This keeps each policy purely optimised for its trade-off, with no
interference from conflicting gradients. Episodes end on acceptance
or when the step limit (\texttt{MAX\_EPISODE\_STEPS}) is reached; a truncated episode counts as non-accepting and receives no special reward. Truncation is reported to the learner as ordinary
termination, a simplification noted in Section~\ref{sec:limitations}.
Per-step logs record reward components, role selections, cost,
duration, and compliance metrics.

A brief terminological note: the \emph{policy} is the fixed set of
network weights produced by training, a function from marking to
action; a \emph{trace} is the sequence of executed events produced by
applying that policy repeatedly within one episode. The same policy
can produce different traces depending on the initial marking and
which actions the shield intercepts along the way. A concrete
worked example of this function evaluated on a real marking is given in Section~\ref{sec:discussion_csp_comparison}.


The Pareto front is extracted from the accepting traces of the final
20\,\% of each run, the converged phase as verified directly against
the per-decile reward statistics of Section~\ref{sec:expense_baseline}
(Appendix~\ref{app:expense_reward_deciles}); truncated episodes are excluded from the Pareto and acceptance statistics. Training budgets per
benchmark are given in Chapter~\ref{chap:experimental_setup}.

% ============================================================
\section{Role-conditioned Extension of \ac{EDCR}}
\label{sec:roles}

This section formalises contribution~C3. Real processes often allow
the same activity to be performed by different resources, a junior
employee or an expert, each with its own cost--duration profile, so
choosing the right resource is part of the optimisation problem. The
environment is therefore extended to joint \texttt{(event, role)}
decisions, a novel extension of \ac{EDCR} graphs in which role
profiles are encoded directly in the process specification while
remaining compatible with existing \ac{EDCR} graphs.

Each role $\rho$ defines a cost multiplier $\mu^\rho_c$ and a
duration multiplier $\mu^\rho_d$:

\begin{equation}
c_{e,\rho} = c_e \cdot \mu^\rho_c, \qquad
d_{e,\rho} = d_e \cdot \mu^\rho_d
\end{equation}

where $c_e$ and $d_e$ are the base values in the graph. As noted in
Section~\ref{sec:reward}, these role-adjusted values enter the reward
before normalisation, so role choice directly scales the per-step
penalty. The observation space is unchanged, the agent sees only the
marking, so role specialisation must be learned from the reward
alone. Table~\ref{tab:role_multipliers} lists the profiles used in
the Loan Application experiment.

\begin{table}[H]
\centering
\caption{Role profiles used in the Loan Application experiment. Both
multipliers are dimensionless ratios applied to the base cost/duration
of Table~\ref{tab:loan_costs}.}
\label{tab:role_multipliers}
{\renewcommand{\arraystretch}{1.5}
\setlength{\tabcolsep}{14pt}
\begin{tabular}{lrrl}
\toprule
\textbf{Role} & \shortstack{\textbf{Cost mult.} \\ (dimensionless)} &
\shortstack{\textbf{Duration mult.} \\ (dimensionless)} & \textbf{Description} \\
\midrule
Junior & $\times 0.7$ & $\times 2.0$ &
    Lower cost, slower execution \\
Expert & $\times 2.0$ & $\times 0.5$ &
    Higher cost, faster execution \\
\bottomrule
\end{tabular}
}
\end{table}


The Junior/Expert structure follows the LoanJS variant of Kirchdorfer
et al.~\cite{KirchdorferLukasFromPolicies}: juniors are cheaper but
slower, seniors faster but more expensive. The multipliers are not
empirically calibrated; they are chosen to create a clear reward
difference between role assignments, enough to test whether the agent
learns meaningful specialisation. Role information is encoded in the
XML definition through the \texttt{roleOptions} extension, which
assigns a cost and duration multiplier to each named role
(Appendix~\ref{app:roleoptions_xml} lists the exact schema). Graphs
without \texttt{roleOptions} fall back to the flat
cost--duration model, preserving compatibility with existing
\ac{EDCR} graphs.


\paragraph{Capacity-constrained allocation.}
With both roles always available, the weight-optimal role of each event
is independent of the others, so role selection scales cost and duration
but does not by itself constitute resource allocation. Contribution~C3
therefore also covers a capacity-constrained variant, in which a
per-case budget on the \emph{Expert} role, enforced by the same shield,
couples the per-event choices into an opportunity-cost allocation
problem. This variant is defined and evaluated in
Section~\ref{sec:loan_dataset} and Chapter~\ref{chap:results}.

\begin{figure}[H]
\centering
\includegraphics[width=0.6\textwidth]{Pictures/thesis_shielding.drawio.png}
\caption{Shielded action execution. When the selected action is
enabled (green path), the engine executes it, updates the marking,
and returns the step reward $r_{\text{step}}$ of
Section~\ref{sec:reward}. When it is not (red path), the shield
intercepts it, the marking is unchanged, and $r = -10$ is returned
with \texttt{done = false}.}
\label{fig:shielding}
\end{figure}
